\section{Background}
\label{sec:background}

\subsection{I/O API of libc}
\label{sec:background:c}

In libc, two different levels of I/O APIs exist:
%
a high-level one using streams of the \C{FILE*} type, and a low-level one using
file descriptors, expressed as integers.
%
This work focuses on the former because it is more widely used due to its
flexibility and convenience, while the latter is recommended for use only when
specifically necessary~\cite{glibc}.

The functions performing stream-based I/O operations are defined in the
\C{stdio.h} header file.
%
The POSIX 2024 standard~\cite{posix2024} specifies 66 API functions to be
declared in this file.
%
However, among them, 6 functions are pure string operations and unrelated to I/O
(e.g., \C{sprintf}),
%
and 3 functions operate on file descriptors rather than streams (e.g.,
\C{dprintf}).
%
For this reason, this work focuses on replacing the remaining 57 functions, such
as \C{fopen}, \C{fread}, \C{fwrite}, and \C{fseek}.
%
The full list is provided in the supplementary material~\cite{supp}.
%
We now describe the types of streams and operations on them.

\paragraph{Origins}

While streams can have various origins, they all have the same type \C{FILE*}.
%
Possible origins include \emph{standard input/output/error}, \emph{files}
(opened by \C{fopen}), and \emph{pipes} (created by \C{popen}).
%
A pipe represents standard input or output of a subprocess, where output is
selected when the second argument to \C{popen} is \CT{"}\C{r}\CT{"}, and input
when it is \CT{"}\C{w}\CT{"}.
%
The code below is well-typed, as streams have the same type:

\begin{lstlisting}[language=C]
FILE *f; f = stdin; f = stdout; f = stderr;
f = fopen(..); f = popen(..);
\end{lstlisting}

\paragraph{Buffering}

Each stream can be either buffered or unbuffered.
%
By default, standard input/output, files, and pipes are buffered, while standard
error is unbuffered.

\paragraph{Capabilities}

Streams can also have various capabilities, but the type does not distinguish
them either.
%
Capabilities include \emph{reading} data from the stream, \emph{writing} data to
the stream, and \emph{seeking} positions in the stream.
%
The following C code is well-typed, although some operations may fail at run
time because the stream does not have the capability:

\begin{lstlisting}[language=C]
FILE *f = ..;
fread(.., f); fwrite(.., f); fseek(f, ..);
\end{lstlisting}

\paragraph{Closing Streams}

Each stream can be closed by passing it to \C{fclose} (for non-pipes) or
\C{pclose} (for pipes).
%
Using a closed stream causes the operation to fail at run time.

\paragraph{Error Checking}

For error checking, each stream internally maintains two indicators:
%
an \emph{end-of-file (EOF) indicator} and an \emph{error indicator}.
%
When an operation on a stream fails, it sets the EOF indicator if the failure
was due to reaching the end of the stream; otherwise, it sets the error
indicator.
%
Calling \C{feof} on a stream checks whether the EOF indicator is set, and
\C{ferror} checks the error indicator.

\subsection{I/O API of Rust std}
\label{sec:background:rust}

\paragraph{Origins}

Rust std provides different types for streams of different origins:
%
\C{Stdin}, \C{Stdout}, and \C{Stderr} for standard input/output/error;
%
\C{File} for files;
%
and \C{ChildStdin} and \C{ChildStdout} for pipes.
%
Assigning a stream to a variable of a different origin results in a type error:

\begin{lstlisting}[language=Rust]
let f: Stdin = stdin(); // ok
f = stdout(); // type error
\end{lstlisting}

\paragraph{Buffering}

Rust std's streams can also be buffered or unbuffered.
%
As in libc, standard input and output are buffered, and standard error is
unbuffered.
%
However, files and pipes are unbuffered.
%
Instead, Rust std provides the \C{BufReader} and \C{BufWriter} types, which wrap
streams with buffers.
%
For example, \C{BufReader}\CT{<}\C{File}\CT{>} is a readable, buffered file.

\paragraph{Capabilities}

Different capabilities are represented with different \emph{traits}, each of
which represents a set of types that provide certain functionalities.
%
Traits resemble type classes in Haskell and serve a similar role to abstract
classes or interfaces in object-oriented languages.
%
In Rust terminology, a type \C{A} is said to \emph{implement} a trait \C{T} if
\C{A} belongs to \C{T}.

To represent capabilities, Rust std provides \C{Read} for reading, \C{Write} for
writing, and \C{Seek} for seeking.
%
\C{Read} is implemented by \C{Stdin}, \C{File}, \C{ChildStdout}, and
\C{BufReader};
%
\C{Write} is implemented by \C{Stdout}, \C{Stderr}, \C{File}, \C{ChildStdin},
and \C{BufWriter}.
%
\C{Seek} is implemented by \C{File}, and also by \C{BufReader} and \C{BufWriter}
if the underlying stream implements \C{Seek}.

In addition to the aforementioned traits, Rust std also provides the \C{BufRead}
trait for buffered reading.
%
It represents the capability of reading the contents of a buffer, which is
useful for implementing operations where the number of bytes to read is
determined by the data itself, e.g., when translating \C{fscanf}.
%
\C{BufRead} is implemented by \C{Stdin} and \C{BufReader}.
%
% \C{BufRead} is implemented by \C{StdinLock}, which can be obtained by locking
% \C{Stdin}, and by \C{BufReader}.

Since a trait is not a type per se, Rust provides two keywords to use a trait as
a type: \C{impl} and \C{dyn}.
%
For a trait $T$, \C{impl} $T$ and \C{dyn} $T$ are types, and $T$ is called
the \emph{bound} of these types.
%
While similar, \C{impl} and \C{dyn} are used for different purposes.

First, \C{impl} can be used only as the type of a function parameter to accept a
value of any type that satisfies the bound.
%
Passing an argument of an incorrect type results in a type error:

\begin{lstlisting}[language=Rust]
fn foo(f: impl Read) { f.read(..); }
foo(stdin()); // ok
foo(stdout()); // type error
\end{lstlisting}

\noindent
%
The use of \C{impl} does not incur any run-time overhead because functions with
\C{impl} are monomorphized at compile time, emitting separate versions for each
type into the binary.

When using \C{impl}, it is also possible to specify multiple bounds for a single
type using the \CT{+} symbol:

\begin{lstlisting}[language=Rust]
fn foo(f: impl Read+Seek) { f.read(..);f.seek(..); }
\end{lstlisting}

\noindent
%
This requires the argument to implement both traits.

On the other hand, \C{dyn} can be used anywhere to define a variable that can
store a value of any type that satisfies the bound.
%
It incurs run-time overhead because the value is associated with a virtual
method table, and utilizing its capabilities (i.e., calling methods) requires
virtual dispatch.
%
Furthermore, since values of different types may have different sizes, a \C{dyn}
type must be wrapped in a pointer type, such as \C{Box}, which represents a
pointer to a heap-allocated object:

\begin{lstlisting}[language=Rust]
let f: Box<dyn Read> = Box::new(stdin()); // ok
f = Box::new(stdout()); // type error
\end{lstlisting}

While \C{dyn} does not allow specifying multiple bounds, we can define a
\emph{subtrait} that has multiple \emph{supertraits} and use this new trait as a
bound.
%
Only types that implement all the supertraits can implement the subtrait, and
the subtrait inherits all the functionalities of the supertraits:

\begin{lstlisting}[language=Rust]
trait ReadSeek : Read + Seek {}
let f:Box<dyn ReadSeek> = ..; f.read(..);f.seek(..);
\end{lstlisting}

\paragraph{Closing Streams}

Rust std does not provide functions to close streams, and each stream is closed
automatically by its destructor when it is \emph{dropped}, i.e., when it is
passed to the \C{drop} function or goes out of scope.
%
Rust's ownership type system prohibits any use of dropped values, preventing the
use of a closed stream at compile time:

\begin{lstlisting}[language=Rust]
let f: File = File::open(..); drop(f);
f.read(..); // type error
\end{lstlisting}

Note that dropping a pointer to a stream, rather than the stream itself, does
not necessarily close it.
%
Raw pointers (unsafe pointers) and references (safe pointers) merely
\emph{borrow} the pointee, and dropping them does not close the stream.
%
However, a \C{Box} \emph{owns} the pointee, and dropping it closes the stream.
%
For this reason, if a variable is only intended to use a stream without closing
it, it should be defined as a borrowing pointer to prevent unintended closure.
%
For example, \C{f1} is a reference in the following code:

\begin{lstlisting}[language=Rust]
let f: File = File::open(..);
let f1: &mut File = &mut f; f1.read(..); drop(f1);
f.read(..); // ok
\end{lstlisting}

\paragraph{Error Checking}

Each operation on a stream returns a \C{Result} value, which is either \C{Ok}
containing the operation's result or \C{Err} containing the reason for the
failure.
%
The \C{Ok} and \C{Err} cases can be distinguished using \emph{pattern matching},
and in the \C{Err} case, inspecting the reason reveals whether the failure was
due to reaching the end of the stream or another issue.
